\section{Waivers, Observability, and Repair Receipts}

This paper reserves \emph{exceptions} for runtime or program errors. Standards deviations are \emph{waivers}: named, scoped, owned, testable, and temporary departures from a required control. A waiver is not a suppression. A suppression hides or mutes a finding. A compensating control accepts a different proof mechanism for the same risk. A waiver records why the normal control cannot be satisfied yet, which compensating control applies, and when the risk must be repaired or renewed.

Waiver lifecycle states are \texttt{proposed}, \texttt{active}, \texttt{renewed}, \texttt{expired}, \texttt{closed}, and \texttt{superseded}. Active waivers require owner approval, scope, reason, proof requirements, residual risk, expiry, and repair plan. Expired waivers fail closed in ratchet and release modes.

Runtime exceptions still matter for repair. HTTP APIs should use RFC 9457 problem details for boundary errors \cite{rfc9457ProblemDetails2023}. Observability should emit OpenTelemetry exception attributes such as exception type, message, stack trace, and error type \cite{otelSemanticConventions2026,otelExceptionAttributes2026}. TypeScript can preserve underlying causes with \texttt{Error.cause} \cite{mdnErrorCause2026}. Rust should prefer typed errors for domain/application code, using context wrappers only at boundaries where opaque provider detail is acceptable \cite{rustByExampleError2026,anyhowDocs2026}.

\begin{table*}[t]
\caption{Minimum repair receipt fields.}
\label{tab:repair-receipts}
\centering
\scriptsize
\begin{tabularx}{\textwidth}{L{0.24\textwidth} Y}
\toprule
\textbf{Field} & \textbf{Purpose} \\
\midrule
\texttt{receipt\_id} & Stable identifier for the repair evidence. \\
\texttt{standard\_version}, \texttt{auditor\_version}, \texttt{schema\_version} & Version bindings for interpretation. \\
\texttt{before\_commit}, \texttt{after\_commit} & Exact repair range. \\
\texttt{started\_at}, \texttt{completed\_at} & Timing evidence and repair half-life input. \\
\texttt{command\_digest}, \texttt{log\_digest}, \texttt{artifact\_digests} & Replay and tamper evidence. \\
\texttt{decision} & \texttt{pass}, \texttt{review}, \texttt{block}, \texttt{ratchet\_fail}, or \texttt{release\_fail}. \\
\texttt{waiver\_closed} & Whether a prior waiver was closed by the repair. \\
\texttt{residual\_risk}, \texttt{redaction\_status} & Remaining risk and evidence privacy state. \\
\texttt{owner\_approval} & Owner route or approval receipt. \\
\texttt{expires\_at} & Required when the repair is partial or waiver-backed. \\
\bottomrule
\end{tabularx}
\end{table*}

Every controlled error that can reach logs, API responses, background jobs, or tests should expose a stable name, stable code, purpose, reason, common fixes, documentation URL, safe user message, correlation ID, and source cause when safe. The point is not verbosity. The point is to make a failing test or production trace directly actionable for the next agent.

Repair is a governed loop, not permission for unattended mutation. \texttt{jankurai repair-plan} reads audit output and turns findings into ordered, schema-validated candidate repairs. \texttt{jankurai repair --dry-run} produces a receipt without touching the repository. Real application is gated behind explicit environment variables, bounded risk, and reviewed flags; git commits, pushes, and draft PR creation require separate opt-ins. Jankurai can prepare a proof-backed branch, commit, and draft PR, but high-risk autonomous mutation and auto-merge remain outside the standard.
